
In this section, we develop estimation and inference results for the closed-form bounds derived in \cref{sec:pid}. For the discrete instrument setting, we leverage DML \citep{chernozhukov2018double} to accommodate high-dimensional covariates, constructing Neyman-orthogonal scores that yield $\sqrt{n}$-consistent and asymptotically normal estimators. In the continuous instrument setting, we characterize the corresponding nonparametric convergence rates. Throughout, we observe an i.i.d. sample $\{O_i\}_{i=1}^n$ of $O=(X,Z,W,Y)\sim\mathbb{P}_{\text{obs}}$, and $\eta$ collects the nuisance functions entering the bound, with $\hat\eta$ its estimate.

\noindent\textit{Remark (notation).} $\eta$ collects the nuisance functions entering the applicable bound and estimation formula, with its components determined by the instrument and outcome regime and $\hat\eta$ denoting their fitted counterparts. In Lemma~\ref{lem:orthogonality_discrete}, Theorem~\ref{thm:dml_normality}: In part~\textup{(i)}, $\eta(x)=\{p(z,x),p_k(x),\pi_k(x),\tau_{j,k}(x),a_{j,k}(z,x),\gamma_{\text{full},k}(x),\gamma_{j,k}(x),\gamma_K(x),\nu_{j,k}(x),M_{j,k}^{\pm}(x),J_{j,k}^{\pm}(x),J_{\text{full},k}^{\pm}(x)\}_{j,k}$; $\pi_k$ and minus-side objects are included only for $k\geqslant1$, and estimated quantiles have $j\geqslant1$, with structural boundary values fixed. In part~\textup{(ii)}, $\eta(x)=\{p(z,x),p_k(x),\pi_k(x),\tau_{j,k}(x),a_{j,k}(z,x),\gamma_{\text{full},k}(x),\gamma_{j,k}(x),\gamma_K(x),P_{1,k}(x)\}_{j,k}$, with $p_0=P_{1,0}=0$ and no fitted $\pi_0$; within-set share directions satisfy $\sum_{z\in T_{j,k}}\delta a_{j,k}(z,x)=0$. In Theorem~\ref{thm:dml_normality}: The unaccented $\eta$ in $\sigma^2=\mathbb E[(\psi(O;\eta)-\underline\theta_{\omega,1})^2]$ denotes $\eta_0$, the true nuisance vector. In Theorem~\ref{thm:continuous_rate}, Corollary~\ref{cor:continuous_primitive_rate}: The nuisance functions are those entering the connected-instrument plug-in representation: $p$, the policy propensity and support endpoints derived from it, $g_1$, and the treated boundary law or its quantile function. $\hat\eta$ is the vector of fitted nuisance functions entering the applicable bound estimator, with its components determined by the instrument and outcome regime. In Theorem~\ref{thm:dml_normality}: The first stage is fitted on $I_1$ and the remaining nuisances are fitted on $I_1$ given the recovered level sets, subject to the stated clipping conventions. Feasible recovery uses $\{X_i:i\in I_2\}$; in the oracle procedure the population level sets are used and $\hat\eta$ is fixed conditional on $I_1$. In Theorem~\ref{thm:continuous_rate}, Corollary~\ref{cor:continuous_primitive_rate}: The fitted nuisances enter (\ref{eq:estimator_continuous_instrument}), including estimates of $p$, the derived policy and support endpoints, $g_1$, and the treated boundary law or its quantile function; they are trained on $I_0\cup I_1$ and held fixed conditional on these training samples for evaluation on $I_2$. In Corollary~\ref{cor:continuous_primitive_rate}: The concrete procedure fits $p$ on $I_0$ and, conditional on that first stage, fits $g_1$ and the treated boundary law on $I_1$, with predictions clipped to their stated supports.


We model the alternative policy as $q(z,x) = \phi(z,x,p(z,x))$, where the known map $\phi:\mathcal Z\times\mathcal X\times[0,1]\to[0,1]$ is jointly measurable. Common examples include propensity shifts $\phi(z,x,p) = p + \alpha$ or $p(1+\alpha)$ on designs for which the shifted propensities remain in $[0,1]$, or setting $\phi(z,x,p) = r(z,x)$ to evaluate a specific alternative policy $r$.

\noindent\textit{Remark (notation).} $\phi:\mathcal Z\times\mathcal X\times[0,1]\to[0,1]$ is the known jointly measurable policy map specifying $q(z,x)=\phi(z,x,p(z,x))$ and the substitution $\hat q(z,x)=\phi(z,x,\hat p(z,x))$. In Lemma~\ref{lem:orthogonality_discrete}, Theorem~\ref{thm:dml_normality}: $\phi(z,x,p)$ is differentiable in $p$, and $\partial\phi/\partial p$ is bounded and uniformly Lipschitz in $p$ over $z\in\mathcal Z$ and $x\in\mathcal X$. In Corollary~\ref{cor:continuous_primitive_rate}: $|\phi(z,x,p')-\phi(z,x,p)|\leqslant L_\phi|p'-p|$ for all $(z,x,p,p')\in\mathcal Z\times\mathcal X\times[0,1]^2$, with $L_\phi<\infty$.


We focus our exposition on the bounds for the PRTE because the derivations for other causal quantities proceed similarly. For simplicity, we first develop theory for the unscaled numerator $\mathbb{E}[Y^{q}-Y]$. In the empirical designs considered below, the PRTE denominator is the known policy-shift size $\alpha > 0$, so the reported bounds and confidence intervals for $\text{PRTE}_\alpha$ are obtained by dividing the corresponding numerator bounds by $\alpha$. More generally, when $\mathbb{E}[q(Z,X)-p(Z,X)]$ is not fixed by design, it is point-identified and can be estimated separately. Thus, our target weight function simplifies to
\begin{align*}
    \omega(x,u) \equiv \mathbb{E}_{\text{obs}}[\indicator(u \leqslant  q(Z,X)) -  \indicator(u \leqslant  p(Z,X)) \mid X=x].
\end{align*}

\subsection{Discrete Instrument Setting}

For a fixed $x$, the weight function $\omega(x,u)$ is a step function that is non-increasing on each latent interval $u \in (p_i(x), p_{i+1}(x)]$, with jumps only at the alternative-policy values $q(z,x)$ falling inside that interval. This piecewise structure allows us to decompose the closed-form lower bound integral into a finite weighted sum over the instrument level sets.

Define $q_{j,i}(x)$ as the $j$-th smallest value of the alternative policy propensity $q(z,x)$ that falls in the baseline bin $[p_i(x), p_{i+1}(x))$ for $i<K$ and in the closed final bin $[p_K(x),1]$ for $i=K$, ordered as follows:
\begin{equation*}
    p_i(x) \leqslant q_{1,i}(x) < \cdots < q_{l_i, i}(x) <  p_{i+1}(x)\quad(i<K),
\end{equation*}
with the boundary conventions $q_{0,i}(x) = p_i(x)$ and $q_{l_i+1,i}(x) = p_{i+1}(x)$. Let the baseline and alternative instrument level sets be
\begin{align*}
    S_k = \{z \in \mathcal{Z} : p(z,x) = p_k(x)\}, \quad T_{j,k} = \{z \in \mathcal{Z} : q(z,x) = q_{j,k}(x)\}, \quad j=1,\ldots,l_k.
\end{align*}
For the final bin, the last displayed strict inequality is replaced by $q_{l_K,K}(x)\leqslant p_{K+1}(x)=1$. Here $q_{0,i}$ and $q_{l_i+1,i}$ are integration boundaries; set $T_{0,i}=T_{l_i+1,i}=\emptyset$. Thus, if $q_{1,i}=p_i$, instruments attaining that value belong only to $T_{1,i}$, while an alternative propensity equal to one belongs to $T_{l_K,K}$ and is retained in the tail term. Values of $u$ exactly on a boundary are Lebesgue-null and do not affect the target integral.

\noindent\textit{Remark (notation).} $q_{j,k}(x)$ denotes a population policy level for $k=0,\ldots,K$ and $j=1,\ldots,l_k$, ordered distinctly within its baseline bin, with $T_{j,k}=\{z\in\mathcal Z:q(z,x)=q_{j,k}(x)\}$; $K,l_k,T_{j,k}$ are invariant in $x$, and $q_{0,k},q_{l_k+1,k}$ are integration boundaries with empty associated instrument sets. In Lemma~\ref{lem:orthogonality_discrete}, Theorem~\ref{thm:dml_normality}: For every nonempty $T_{j,k}$ in a nonfinal bin $k<K$, $q_{j,k}(x)\neq p_k(x)$; the structural boundary $q_{0,k}=p_k$ is therefore distinct from every regular policy boundary used by the displayed scores. $l_k\in\mathbb N_0$, for $k=0,\ldots,K$, counts distinct population policy levels rather than instruments, is invariant in $x$ under the stipulated invariant ordering, and excludes structural integration boundaries while retaining an actual policy value equal to one in the final bin. $S_k=\{z\in\mathcal Z:p(z,x)=p_k(x)\}$ denotes the population baseline instrument level set for $k=1,\ldots,K$, invariant in $x$ under the maintained equality-class convention; $S_0=\{z\in\mathcal Z:p(z,x)=0\}$, when nonempty, is the structural-zero class, and $\widehat S_k$ denotes a recovered set rather than $S_k$ itself. In Lemma~\ref{lem:orthogonality_discrete}, Theorem~\ref{thm:dml_normality}: At $k=0$, minus-side conditional expectations and residuals are omitted and assigned zero only for notation; no treated conditional regression on $S_0$ or division by $\pi_0(X)$ is used. $T_{j,k}$ denotes the population alternative-policy equality class for $k=0,\ldots,K$ and $j=1,\ldots,l_k$, invariant in $x$ under the stipulated invariant ordering; structural integration-boundary sets are empty, an actual policy value equal to $p_k(x)$ belongs to $T_{1,k}$, and an actual value equal to one is retained in $T_{l_K,K}$. In Lemma~\ref{lem:orthogonality_discrete}, Theorem~\ref{thm:dml_normality}: For $k=0,\ldots,K-1$, every nonempty $T_{j,k}$ has $q_{j,k}(x)\neq p_k(x)$ for all $x\in\mathcal X$. $S$ is an auxiliary $\mathrm{Unif}(0,1)$ draw independent of $U\sim\nu_{k\mid x}$ for each fixed $x$ and nonempty gap $G_k(x)$, used to randomize ties in $R_{k,x}(U,S)$ and integrated out in the attaining outcome kernel.


The DML scores developed after the reduction differentiate $q(z,x)=\phi(z,x,p(z,x))$ with respect to the propensity score, so we impose the following condition only for inference.

\begin{assumption}[Policy-Map Smoothness]\label{asp:phi_smooth}
    The map $\phi(z,x,p)$ is differentiable with respect to $p$, and $\frac{\partial \phi}{\partial p}$ is bounded and uniformly Lipschitz in $p$ for all $z \in \mathcal{Z}$ and $x \in \mathcal{X}$.
\end{assumption}

\noindent\textit{Remark.} For every $z\in\mathcal Z$ and $x\in\mathcal X$, $p\mapsto\phi(z,x,p)$ is differentiable on $(0,1)$ and has finite right and left derivatives at $p=0$ and $p=1$, respectively; these are its endpoint values of $\partial\phi/\partial p$. There exist finite constants $C,L\geqslant0$ such that, for all $z\in\mathcal Z$, $x\in\mathcal X$, and $p,p'\in[0,1]$, $|\partial\phi/\partial p(z,x,p)|\leqslant C$ and $|\partial\phi/\partial p(z,x,p)-\partial\phi/\partial p(z,x,p')|\leqslant L|p-p'|$.

This condition controls the pathwise derivative and its second-order remainder.

We impose the following separation condition to consistently estimate the level sets.

\begin{assumption}[Propensity-Level Separation]\label{asp:prop_gap}
    \begin{enumerate}[label=(\roman*)]
        \item There exists a constant $c_{\text{gap}} > 0$ such that, for all $x \in \mathcal{X}$ and all $v,v' \in R_x = \{p(z,x),q(z,x)\}_{z \in \mathcal{Z}}$, either $v'=v$ or $|v-v'|>c_{\text{gap}}$. Moreover, the equality classes and their strict ordering are invariant in $x$: the sets $S_k$, $T_{j,k}$ and their indices $K,l_k$ do not depend on $x$.
        \item For every $k=0,\ldots,K-1$ and every nonempty $T_{j,k}$, $q_{j,k}(x)\neq p_k(x)$ for all $x$. Under part~\textup{(i)}, this implies $q_{j,k}(x)>p_k(x)+c_{\text{gap}}$ for $k\geqslant1$.
    \end{enumerate}
\end{assumption}

\noindent\textit{Remark (notation).} $R_x=\{p(z,x),q(z,x):z\in\mathcal Z\}$ is the set of distinct baseline and alternative-policy propensity values at the covariate value $x$. In Proposition~\ref{prop:discrete_reduction}, Lemma~\ref{lem:orthogonality_discrete}, Theorem~\ref{thm:dml_normality}: Distinct elements of $R_x$ are separated by more than $c_{\text{gap}}$, uniformly in $x$; the equality classes and their strict ordering are invariant in $x$, although the numerical propensity values may vary with $x$. $c_{\text{gap}}$ is one common positive separation constant for distinct values in the combined baseline and alternative-policy range, uniformly over $x\in\mathcal X$. In Theorem~\ref{thm:dml_normality}: Level recovery joins averaged fitted values lying within $c_{\mathrm{gap}}/2$ of each other, where $c_{\mathrm{gap}}$ denotes this same separation constant.


Part~\textup{(i)} implies the stable level sets needed by the finite-dimensional reduction and the clustering step below. We therefore suppress the dependence on $x$ in $S_k$, $T_{j,k}$, $K$, and $l_k$ from now on. Part~\textup{(ii)} keeps a nonempty pooled policy boundary distinct from the structural boundary $q_{0,k}=p_k$ in the displayed continuous-outcome quantile score and binary active-set score. We impose this condition only to simplify the clustering labels and score formulas; it is not required for identification. It can be removed by using a more elaborate clustering method that detects the coinciding boundaries.

Because $\mathcal{Z}$ is finite, \cref{asp:continuity_p} holds trivially and every conditional LIV interval in \cref{thm:covariate_continuous_lower_bound} is degenerate, so that theorem specializes to the covariate version of \cref{thm:discrete_lower_bound}; we refer to it in this form throughout the subsection. The resulting bound admits a representation as an expectation of a finite weighted sum over instrument level sets and their sub-intervals. Before stating this representation, we introduce the nuisance components that carry the observable content of the bound.

\paragraph{Nuisance components.} The representation below is assembled from five objects, defined for each interval index $k \in \{0,\dots,K\}$ and sub-interval index $j \in \{0,\dots,l_k\}$. The first two are the aggregate weights that the step weight $\omega$ assigns to instruments,
\begin{align*}
    \gamma_{\text{full},k}(x) &\coloneqq \mathbb{P}_{\text{obs}}\!\left(Z \in \bigcup_{i=k+1}^K \bigcup_{l=0}^{l_i} T_{l,i} \mathrel{\Big|} X=x\right) - \mathbb{P}_{\text{obs}}\!\left(Z \in \bigcup_{i=k+1}^K S_i \mathrel{\Big|} X=x\right), \\
    \gamma_{j,k}(x) &\coloneqq \mathbb{P}_{\text{obs}}\!\left(Z \in \bigcup_{l=j}^{l_k} T_{l,k} \mathrel{\Big|} X=x\right),
\end{align*}
where $\gamma_{\mathrm{full},k}\in[-1,1]$ because it differences two conditional probabilities. The next two carry the outcome distribution: with the relative threshold $\kappa_{j,k}(x) \coloneqq (q_{j,k}(x) - p_k(x))/(p_{k+1}(x) - p_k(x)) \in [0,1]$,
\begin{align*}
    J_{\text{full},k}(x) &\coloneqq (p_{k+1}(x)-p_k(x))\,\mathbb{E}_{\mu_{1,k\mid x}}[Y], \\
    J_{j,k}(x) &\coloneqq (p_{k+1}(x)-p_k(x))\int_{\kappa_{j,k}(x)}^{\kappa_{j+1,k}(x)} Q_{Y,k\mid x}(u)\,\mathrm{d}u,
\end{align*}
where $\mu_{1,k\mid x}$ is the identified conditional complier distribution of (\ref{eq:gap_identified_measure_cov}) and $Q_{Y,k\mid x}$ is its quantile function. The fifth is the trivial-bound boundary term
\begin{equation*}
    \Delta_K(x) \coloneqq y_{\min}\, \mathbb{E}_{\text{obs}}\!\left[ \indicator\!\Big(Z \in \bigcup_{j=1}^{l_K} T_{j,K}\Big) (q(Z,X) - p_K(X)) \mathrel{\Big|} X=x \right].
\end{equation*}
This is the tail term of \cref{thm:covariate_continuous_lower_bound} under the PRTE weight: since $p(z,x)\leqslant p_K(x)$ for every $z$, the weight satisfies $\omega(x,u)=\mathbb{E}_{\text{obs}}[\indicator(u\leqslant q(Z,X))\mid X=x]\geqslant0$ on $(p_K(x),1]$, so the $y_{\max}$ branch of the tail integral is inactive, and Fubini converts $y_{\min}\int_{p_K(x)}^1\omega(x,u)\,\mathrm{d}u$ into the displayed expectation.

\noindent\textit{Remark (notation).} $\gamma_{\mathrm{full},k}(x)$ is the population difference of the two conditional probabilities in its defining display, taking values in $[-1,1]$ for $k=0,\ldots,K$, with $\gamma_{\mathrm{full},K}\equiv0$ by empty unions; the full-interval products use $k=0,\ldots,K-1$. $\kappa_{j,k}(x)$ is the population relative threshold for a nonempty baseline interval, indexed by $k=0,\ldots,K-1$ and $j=0,\ldots,l_k+1$, with $\kappa_{0,k}(x)=0$ and $\kappa_{l_k+1,k}(x)=1$. $J_{\text{full},k}(x)$ is the population full-interval contribution $(p_{k+1}(x)-p_k(x))\mathbb E_{\mu_{1,k\mid x}}[Y]$ for $k=0,\ldots,K-1$. In Proposition~\ref{prop:J_specialization}, Lemma~\ref{lem:orthogonality_discrete}, Theorem~\ref{thm:dml_normality}: In the binary-outcome case, $J_{\text{full},k}(x)=P_{1,k+1}(x)-P_{1,k}(x)$, where $P_{1,k}(x)=\mathbb E_{\text{obs}}[YW\mid Z\in S_k,X=x]$ for $k\geqslant1$ and $P_{1,0}(x)=0$. $\Delta_K(x)=y_{\min}\mathbb E_{\text{obs}}[\max\{0,q(Z,X)-p_K(X)\}\mid X=x]$ is the conditional lower-bound contribution on $(p_K(x),1]$, with the final policy bin $[p_K(x),1]$ retaining instruments whose alternative propensity equals one.


Estimating these five requires three further level-set quantities that do not appear in the representation itself---the tail weight $\gamma_K(x)\coloneqq\mathbb{P}_{\text{obs}}(Z\in\bigcup_{j=1}^{l_K}T_{j,K}\mid X=x)$, the within-set assignment shares $a_{j,k}(z,x)\coloneqq\mathbb{P}_{\text{obs}}(Z=z\mid Z\in T_{j,k},X=x)$, and the pooled representation $q_{j,k}(x)=\sum_{z\in T_{j,k}}a_{j,k}(z,x)q(z,x)$ that extends the policy level regularly when estimated instrument-specific levels are not exactly tied. The latter two are detailed in \cref{app:procedure_continuous_outcome}, and all three enter only through the scores.

\noindent\textit{Remark (notation).} $\gamma_K(x)$ is the population conditional probability of $\bigcup_{j=1}^{l_K}T_{j,K}$ in the final closed policy bin $[p_K(x),1]$, taking values in $[0,1]$ and equaling zero when $l_K=0$; actual policy levels equal to one are included. $a_{j,k}(z,x)=\mathbb P_{\mathrm{obs}}(Z=z\mid Z\in T_{j,k},X=x)$ is the population assignment-share function for each nonempty policy set, indexed by $k=0,\ldots,K$ and $j=1,\ldots,l_k$, with $a_{j,k}(z,x)=0$ outside $T_{j,k}$ and $\sum_{z\in T_{j,k}}a_{j,k}(z,x)=1$ wherever the conditioning probability is positive.


\begin{proposition}[Discrete-Instrument Reduction to Nuisance Functional]\label{prop:discrete_reduction}
    Suppose \cref{asp:structure} and part~\textup{(i)} of \cref{asp:prop_gap} hold and $\mathcal{Z}$ is finite. Then the sharp lower bound in \cref{thm:covariate_continuous_lower_bound} admits the representation
    \begin{align}\label{eq:target_discrete_iv}
        \underline{\theta}_{\omega,1} = \mathbb{E}\Bigg[\sum_{k=0}^{K-1}\bigg(\gamma_{\text{full},k}(X)\, J_{\text{full},k}(X) + \sum_{j=1}^{l_k} \gamma_{j,k}(X)\, J_{j-1,k}(X)\bigg) + \Delta_K(X) \Bigg].
    \end{align}
    The sharp upper bound admits an analogous representation, obtained by replacing $Q_{Y,k\mid x}(u)$ by $Q_{Y,k\mid x}(1-u)$ inside $J_{j,k}$ and swapping $y_{\min}$ with $y_{\max}$ in $\Delta_K$.
\end{proposition}

\noindent\textit{Remark.} The identified outcome components $J_{\text{full},k}$ and $J_{j,k}$ are defined for $0\leqslant k<K$, with $0\leqslant j\leqslant l_k$ for $J_{j,k}$; the final interval is represented by $\Delta_K$, without an identified law $\mu_{1,K\mid x}$.

\noindent\textit{Remark (notation).} $\gamma_{j,k}(x)$ is the population conditional probability of $\bigcup_{l=j}^{l_k}T_{l,k}$, taking values in $[0,1]$ for $k=0,\ldots,K$ and $j=0,\ldots,l_k$, with $T_{0,k}=\emptyset$; the sub-interval products use $k=0,\ldots,K-1$ and $j=1,\ldots,l_k$. $J_{j,k}(x)=(p_{k+1}(x)-p_k(x))\int_{\kappa_{j,k}(x)}^{\kappa_{j+1,k}(x)}Q_{Y,k\mid x}(u)\,\mathrm du$ is the population lower-bound subinterval contribution on a nondegenerate identified bin, with structural thresholds $\kappa_{0,k}=0$ and $\kappa_{l_k+1,k}=1$. In Proposition~\ref{prop:discrete_reduction}, Proposition~\ref{prop:J_specialization}, Lemma~\ref{lem:orthogonality_discrete}, Theorem~\ref{thm:dml_normality}: The bound uses $J_{j-1,k}$ for $0\leq k<K$ and $1\leq j\leq l_k$. In the continuous-outcome representation, for $0\leq j<l_k$ and under the stated continuity conditions, $J_{j,k}=(p_{k+1}-p_k)\mathbb E_{\mu_{1,k\mid X}}[Y\indicator(\nu_{j,k}<Y\leqslant\nu_{j+1,k})]$; in the binary-outcome case, $J_{j-1,k}=h_{j,k}^{-}$. In Lemma~\ref{lem:orthogonality_discrete}, Theorem~\ref{thm:dml_normality}: In the continuous-outcome score construction the plug-in component is $J_{j,k}=p_{k+1}J_{j,k}^{+}-p_kJ_{j,k}^{-}$; at $k=0$, the minus-side conditional expectation is omitted and set to zero for notation.


\cref{prop:discrete_reduction} is the key analytical step that converts the optimal-transport integral in \cref{thm:covariate_continuous_lower_bound} into a Neyman-orthogonalizable functional of standard nuisance quantities---conditional probabilities, conditional means, and conditional quantile integrals. The proof, given in \cref{app:closed_form_rewriting} of \oalocation, exploits the piecewise-constant, monotone structure of $\omega(X,\cdot)$ on each baseline interval and identifies the step-function value attached to each sub-interval $(q_{j,k}(X), q_{j+1,k}(X))$.

The sub-interval term $J_{j,k}$ still involves the conditional quantile function $Q_{Y,k\mid X}$, which is delicate to estimate directly. The next proposition replaces it with conditional expectations in the two outcome regimes of interest.

\begin{proposition}[Specialization of the Sub-Interval Contribution]\label{prop:J_specialization}
    Fix $k \in \{0,\dots,K-1\}$. Set the lower boundary $\nu_{0,k}(x)\equiv y_{\min}$, and let $\nu_{j,k}(x) \coloneqq Q_{Y,k\mid x}(\kappa_{j,k}(x))$ for $j \in \{1,\dots,l_k\}$.
    \begin{enumerate}[label=(\roman*)]
        \item \textbf{Continuous outcome.} If $\mu_{1,k\mid X}$ is continuous at $\nu_{j,k}(X)$ and $\nu_{j+1,k}(X)$ almost surely, then, for $j \in \{0,\dots,l_k-1\}$,
        \begin{equation}\label{eq:J_cexpectation}
            J_{j,k}(X) = (p_{k+1}(X) - p_k(X))\,\mathbb{E}_{\mu_{1,k\mid X}}\!\big[ Y\,\indicator\!\big(\nu_{j,k}(X) < Y \leqslant \nu_{j+1,k}(X)\big) \big].
        \end{equation}
        \item \textbf{Binary outcome.} If $Y \in \{0,1\}$, then $J_{\text{full},k}(X) = P_{1,k+1}(X) - P_{1,k}(X)$ with $P_{1,k}(x) \coloneqq \mathbb{E}_{\text{obs}}[YW \mid Z \in S_k, X=x]$ for $k\geqslant1$ and the boundary convention $P_{1,0}(x)\equiv0$, and for $j \in \{1,\dots,l_k\}$,
        \begin{equation}\label{eq:h_jk_binary}
            J_{j-1,k}(X) = h_{j,k}^{-}(X) \coloneqq \max\!\Big(0,\; q_{j,k}(X) - \max\!\big(q_{j-1,k}(X),\; p_{k+1}(X) - J_{\text{full},k}(X)\big)\Big).
        \end{equation}
        Consequently, (\ref{eq:target_discrete_iv}) collapses to
        \begin{align}\label{eq:binary_lower_bound}
            \underline{\theta}_{\omega,1} = \mathbb{E}\Bigg[\sum_{k=0}^{K-1}\bigg(\gamma_{\text{full},k}(X)\, J_{\text{full},k}(X) + \sum_{j=1}^{l_k} \gamma_{j,k}(X)\, h_{j,k}^{-}(X)\bigg) + \Delta_K(X)\Bigg].
        \end{align}
    \end{enumerate}
\end{proposition}

\noindent\textit{Remark.} The boundary quantile uses the bounded-domain lower generalized inverse $Q_{Y,k\mid x}(t)=\inf\{y\in\mathcal Y:\mu_{1,k\mid x}(\mathcal Y\cap(-\infty,y])\geqslant t\}$ for $0\leqslant t<1$, so $Q_{Y,k\mid x}(0)=y_{\min}$. For each fixed $k\in\{0,\dots,K-1\}$ and $j\in\{0,\dots,l_k-1\}$ to which part~(i) is applied, its continuity hypothesis is $\mu_{1,k\mid X}(\{\nu_{j,k}(X)\})=\mu_{1,k\mid X}(\{\nu_{j+1,k}(X)\})=0$ almost surely under the marginal law of $X$. For $0<t<1$, use the standard lower generalized quantile $Q_{Y,k\mid x}(t)=\inf\{y\in\mathbb R:\mu_{1,k\mid x}(\mathcal Y\cap(-\infty,y])\geqslant t\}$. For $j=1,\ldots,l_k$, interpret $\nu_{j,k}(x)=Q_{Y,k\mid x}(\kappa_{j,k}(x))$ using $Q_{Y,k\mid x}(t)=\inf\{y\in\mathbb R:\mu_{1,k\mid x}(\mathcal Y\cap(-\infty,y])\geqslant t\}$ for $0<t<1$, with $\nu_{0,k}(x)=y_{\min}$.

\noindent\textit{Remark (notation).} $\nu_{j,k}$ denotes the population boundary function for $k=0,\ldots,K-1$ and $j=0,\ldots,l_k$, with $\nu_{0,k}(x)\equiv y_{\min}$ and $\nu_{j,k}(x)=Q_{Y,k\mid x}(\kappa_{j,k}(x))$ for $j\geqslant1$; $\hat\nu_{j,k}$ denotes its estimate. In Proposition~\ref{prop:J_specialization}: In part (i), the gap law $\mu_{1,k\mid X}$ is continuous at the two boundaries $\nu_{j,k}(X)$ and $\nu_{j+1,k}(X)$ used in each band, almost surely. In Lemma~\ref{lem:orthogonality_discrete}, Theorem~\ref{thm:dml_normality}: In the continuous-outcome part, the relevant treated conditional laws have uniformly bounded densities in fixed neighborhoods of $\nu_{j,k}(X)$ for $j=1,\ldots,l_k$, and no atom at the fixed boundary $\nu_{0,k}=y_{\min}$; only the boundaries with $j\geqslant1$ are estimated. For $k=1,\ldots,K$, $P_{1,k}(x)=\mathbb E_{\mathrm{obs}}[YW\mid Z\in S_k,X=x]$ is the population conditional treated-success regression in the binary-outcome setting, while $P_{1,0}(x)\equiv0$ is a structural boundary convention rather than a conditional regression on $S_0$. $h_{j,k}^{-}(x)$ is the binary-outcome population subinterval contribution for $k=0,\ldots,K-1$ and $j=1,\ldots,l_k$, with $q_{0,k}(x)=p_k(x)$, and equals $J_{j-1,k}(x)$.


Part~(i) converts the sub-interval quantile integral into a plain conditional expectation of $Y$ on the quantile-defined band $(\nu_{j,k}, \nu_{j+1,k}]$, so only the $l_k+1$ boundary quantiles $\nu_{j,k}$---rather than the full quantile process---enter the estimation. Part~(ii) leverages the step-function form of $Q_{Y,k\mid X}$ when $Y$ is binary: the complier mean $\mathbb{E}_{\mu_{1,k\mid X}}[Y]$ reduces to a difference of observed joint probabilities, and the integral of the step function over $[\kappa_{j-1,k}, \kappa_{j,k}]$ evaluates to the closed-form $h_{j,k}^{-}$. Both parts are proved in \cref{app:J_specialization} of \oalocation.

Together, \cref{prop:discrete_reduction,prop:J_specialization} expose (\ref{eq:target_discrete_iv}) as a functional of conditional expectations, boundary quantiles, and instrument-level-set probabilities---the form required for DML-style $\sqrt{n}$-inference. In both outcome regimes the estimator takes the same aggregated Neyman-orthogonal form,
\begin{equation}\label{eq:estimator_dml}
    \hat{\underline{\theta}}_{\omega, 1} = \frac{1}{|I_2|} \sum_{i \in I_2} \Bigg[ \sum_{k=0}^{K-1}\bigg(\hat{\psi}_{\text{full},k}^{\text{prod}}(O_i; \hat{\eta}) + \sum_{j=1}^{l_k} \hat{\psi}_{j,k}^{\text{prod}}(O_i; \hat{\eta})\bigg) + \hat{\psi}_{\Delta_K}(O_i; \hat{\eta}) \Bigg],
\end{equation}
where each $\hat{\psi}$ is a Neyman-orthogonal score that corrects for the estimation error of its associated nuisance. What differs across the two regimes is the nuisance list $\hat\eta$, and hence the score formulas: these are given in \cref{app:scores_continuous} for a continuous outcome and in \cref{app:scores_discrete} for a discrete one. We treat the two regimes in turn and then state a single asymptotic normality result covering both.

\noindent\textit{Remark (notation).} $\hat{\underline\theta}_{\omega,1}$ denotes the sample-based estimator of the treated-component sharp lower bound, formed by averaging fitted score contributions over $I_2$. In Theorem~\ref{thm:dml_normality}: $\hat{\underline\theta}_{\omega,1}=|I_2|^{-1}\sum_{i\in I_2}\left[\sum_{k=0}^{K-1}\left(\hat\psi_{\text{full},k}^{\text{prod}}(O_i;\hat\eta)+\sum_{j=1}^{l_k}\hat\psi_{j,k}^{\text{prod}}(O_i;\hat\eta)\right)+\hat\psi_{\Delta_K}(O_i;\hat\eta)\right]$, with the nuisance list and score formulas appropriate to the continuous or binary outcome regime. In Theorem~\ref{thm:continuous_rate}, Corollary~\ref{cor:continuous_primitive_rate}: $\hat{\underline\theta}_{\omega,1}=|I_2|^{-1}\sum_{i\in I_2}(\hat\psi_1(O_i)+\hat\psi_2(O_i)+\hat\psi_3(O_i))$, where the three contributions substitute estimated nuisances into the three terms of equation~\ref{eq:continuous_simplified_target}, with predictions clipped to their known supports. In Corollary~\ref{cor:continuous_primitive_rate}: The nuisance estimates are those supplied by Algorithm~\ref{alg:continuous_instrument}, including the local-linear regression and treated boundary conditional quantile construction. $n$ is the total observed sample size indexing the estimation sequences and their asymptotics. In Theorem~\ref{thm:dml_normality}: $n$ ranges over sample sizes permitting the stipulated equal-half split, so $n$ is even. $\hat\psi$ denotes the total fitted integrand averaged over $I_2$ to form $\hat{\underline\theta}_{\omega,1}$. In Theorem~\ref{thm:dml_normality}: $\hat\psi(O;\hat\eta)=\sum_{k=0}^{K-1}\bigl(\hat\psi_{\mathrm{full},k}^{\mathrm{prod}}(O;\hat\eta)+\sum_{j=1}^{l_k}\hat\psi_{j,k}^{\mathrm{prod}}(O;\hat\eta)\bigr)+\hat\psi_{\Delta_K}(O;\hat\eta)$, using the continuous-outcome or binary-outcome score formulas according to the outcome regime. In the oracle procedure used in the proof, this fitted function is fixed conditional on $I_1$; the feasible procedure uses recovered level sets. In Theorem~\ref{thm:continuous_rate}: $\hat\psi=\hat\psi_1+\hat\psi_2+\hat\psi_3$, obtained by substituting fitted nuisances into the three terms of (\ref{eq:continuous_simplified_target}); this is the plug-in integrand, fixed conditional on $I_0,I_1$, rather than the finite-instrument orthogonal score.


\subsubsection{Continuous Outcome}\label{subsubsec:continuous_outcome}

We first treat the case where the conditional complier distribution $\mu_{1,k\mid X}$ is continuous at the boundary quantiles $\nu_{j,k}(X)$, so \cref{prop:J_specialization}(i) applies. Recall that $\nu_{j,k}(x)=Q_{Y,k\mid x}(\kappa_{j,k}(x))$ denotes the boundary quantile introduced in \cref{prop:J_specialization}. The representation~(\ref{eq:J_cexpectation}) converts each sub-interval term $J_{j,k}(X)$ into a conditional expectation over the band $(\nu_{j,k}(X), \nu_{j+1,k}(X)]$, thereby reducing the nuisance list from the full quantile process to the boundary quantiles $\nu_{j,k}$ alone. The Neyman-orthogonal score construction additionally requires the conditional indicator expectations
\begin{align*}
    M_{j,k}^{+}(x) &\coloneqq \mathbb{E}[\indicator(Y \leqslant \nu_{j,k}(x)) \mid Z \in S_{k+1}, W=1, X=x], \\
    M_{j,k}^{-}(x) &\coloneqq \mathbb{E}[\indicator(Y \leqslant \nu_{j,k}(x)) \mid Z \in S_k, W=1, X=x],
\end{align*}
which enter the Riesz representers $\alpha_{p_k}, \alpha_{p_{k+1}}$ for the pathwise derivatives with respect to the interval boundaries $p_k, p_{k+1}$, and whose own estimation error is corrected by the residuals $IF_{M^{\pm}_{j,k}}$ (see \cref{app:scores_continuous}).

\noindent\textit{Remark (notation).} $M_{j,k}^{+}$ is the population conditional indicator-mean function at the true boundary $\nu_{j,k}$, indexed by $k=0,\ldots,K-1$ and $j=0,\ldots,l_k$, with conditioning on $Z\in S_{k+1}$, $W=1$, and $X=x$. In Lemma~\ref{lem:orthogonality_discrete}, Theorem~\ref{thm:dml_normality}: In the continuous-outcome part, the no-atom condition at $y_{\min}$ and $\nu_{0,k}=y_{\min}$ imply $M_{0,k}^{+}(X)=0$ almost surely. $M_{j,k}^{-}(x)$ is the population conditional indicator mean among treated observations with $Z\in S_k$ at the continuous-outcome boundary $\nu_{j,k}(x)$ for $k\geqslant1$; at $k=0$ this minus-side regression is omitted and defined as zero solely for notational convenience. $\alpha_{p_k}$ is the population lower-boundary representer for the particular component $J_{j,k}$, with the dependence on $j$ suppressed in its notation, for $k=1,\ldots,K-1$ and $j=0,\ldots,l_k-1$; there is no $p_0$ perturbation or correction. $\alpha_{p_{k+1}}$ denotes the upper interval-boundary representer for the continuous-outcome component $J_{j,k}$, with suppressed dependence on $j$, indexed by $k=0,\ldots,K-1$ and $j=0,\ldots,l_k-1$; it remains an upper-side representer at $k=0$.


The functional form in~(\ref{eq:target_discrete_iv}) and~(\ref{eq:J_cexpectation}) fixes the nuisance list $\hat\eta$ entering~(\ref{eq:estimator_dml}); the rest of this subsection develops the corresponding estimators, trained on an auxiliary sample $I_1$.

\paragraph{Estimation Procedure.} Following the standard DML template, we randomly partition the sample into an auxiliary set $I_1$ and a main estimation set $I_2$. On $I_1$ we estimate the propensity score $p$, clip it to $[0,1]$, and set $\hat q=\phi(\cdot,\cdot,\hat p)$; using only the covariates in $I_2$, we average $\hat p$ and $\hat q$ over $X$ and recover the level sets $\widehat S_k$ and $\widehat T_{j,k}$ as connected components of the graph that joins two averaged values lying within $c_{\mathrm{gap}}/2$ of each other. Given those sets we fit the remaining nuisances on $I_1$---the interval boundaries $\hat p_k$, the probability weights $\hat\gamma$, $\hat\pi_k$, $\hat\tau_{j,k}$ and within-set shares $\hat a_{j,k}$, the boundary quantiles $\hat\nu_{j,k}$, and the conditional expectations $\hat M^{\pm}_{j,k}$, $\hat J^{\pm}_{j,k}$, $\hat J^{\pm}_{\mathrm{full},k}$---clipping each to its known range, and finally average the plug-in orthogonal scores over $I_2$ to form~(\ref{eq:estimator_dml}). \cref{alg:dml_continuous_outcome} states the procedure step by step; the explicit graph construction, the moment equation defining $\nu_{j,k}$, the conditional-expectation components, and the Neyman-orthogonal scores $\psi^{\mathrm{prod}}$ are given in \cref{app:procedure_continuous_outcome} and \cref{app:scores_continuous}.

\noindent\textit{Remark (notation).} $\psi$ denotes the uncentered total population integrand associated with the estimator of $\underline\theta_{\omega,1}$, with its nuisance arguments supplied explicitly when needed. In Lemma~\ref{lem:orthogonality_discrete}, Theorem~\ref{thm:dml_normality}: $\psi(O;\eta)=\sum_{k=0}^{K-1}\bigl(\psi_{\mathrm{full},k}^{\mathrm{prod}}(O;\eta)+\sum_{j=1}^{l_k}\psi_{j,k}^{\mathrm{prod}}(O;\eta)\bigr)+\psi_{\Delta_K}(O;\eta)$, with the component formulas in Section~\ref{app:scores_continuous} for continuous outcomes and Section~\ref{app:scores_discrete} for binary outcomes; $\eta_0$ supplies the true nuisance values. In Theorem~\ref{thm:continuous_rate}, Corollary~\ref{cor:continuous_primitive_rate}: $\psi(O)=-\int_{\min\{\underline p(X),q(Z,X)\}}^{\underline p(X)}Q_{1,\underline p\mid X}(v/\underline p(X))\,\mathrm dv+g_1(q^{\mathrm{tr}}(Z,X),X)-g_1(P,X)+y_{\min}(q(Z,X)-\overline p(X))_+$, with the first integral zero when $\underline p(X)=0$; this is the population plug-in integrand. $\hat p:\mathcal Z\times\mathcal X\to[0,1]$ is a fitted estimator of the conditional treatment propensity $p(z,x)$, with fitted values clipped to $[0,1]$. In Lemma~\ref{lem:orthogonality_discrete}, Theorem~\ref{thm:dml_normality}: $\hat p$ is the pre-recovery first-stage fit trained on $I_1$; the level-recovery averages evaluate this fit at the covariates in $I_2$. In Theorem~\ref{thm:dml_normality}: $\max_{z\in\mathcal Z}\|\hat p(z,\cdot)-p(z,\cdot)\|_{L_2}=o_P(n^{-1/4})$. In Theorem~\ref{thm:continuous_rate}, Corollary~\ref{cor:continuous_primitive_rate}: $\hat p$ is trained on the propensity sample $I_0$ and is held fixed conditional on $I_0$ when constructing the generated propensity index and localization rules on $I_1$. In Corollary~\ref{cor:continuous_primitive_rate}: $\|\hat p-p\|_\infty=O_P(r_{p,n})$. $\hat q(z,x)=\phi(z,x,\hat p(z,x))\in[0,1]$ is the fitted alternative propensity obtained by substituting the clipped baseline propensity estimator into the known policy map. In Lemma~\ref{lem:orthogonality_discrete}, Theorem~\ref{thm:dml_normality}: $\hat q$ is formed from the pre-recovery propensity fit $\hat p$ trained on $I_1$; the level-recovery average evaluates this fitted function at the covariates in $I_2$. In Theorem~\ref{thm:continuous_rate}, Corollary~\ref{cor:continuous_primitive_rate}: $\hat q$ is formed from $\hat p$ trained on $I_0$; the notation $\hat q=\phi(Z,X,\hat p(Z,X))$ denotes evaluation of this fitted function at $(Z,X)$. $\hat\pi_k(x)$ estimates $\mathbb P_{\mathrm{obs}}(Z\in S_k\mid X=x)$ by fitting the indicator of the recovered set $\widehat S_k$ on $I_1$, for $k=1,\ldots,K$, and is clipped first to $[0,1]$ and then to $[c_\pi/2,1]$; no denominator estimator $\hat\pi_0$ is required at the structural boundary. $\hat\tau_{j,k}(x)$ estimates $\mathbb P_{\mathrm{obs}}(Z\in T_{j,k}\mid X=x)$ by fitting the indicator of the recovered set $\widehat T_{j,k}$ on $I_1$, for nonempty policy level sets with $k=0,\ldots,K$ and $j=1,\ldots,l_k$, and is clipped first to $[0,1]$ and then to $[c_\pi/2,1]$. $\widehat S_k$ is the sample-dependent baseline level-set estimate obtained from the pre-recovery fits on $I_1$ and their averages over the covariates in $I_2$, with a threshold-defined structural-zero set and positively indexed baseline components ordered by their averages. $\widehat T_{j,k}$ is the sample-dependent policy level-set estimate obtained from the pre-recovery fits on $I_1$ and their averages over the covariates in $I_2$; it depends on both of these inputs and is indexed by the recovered component ordering and baseline bins.


Let $\pi_k(x) \coloneqq \mathbb{P}_{\text{obs}}(Z \in S_k \mid X=x)$ and $\tau_{j,k}(x) \coloneqq \mathbb{P}_{\text{obs}}(Z \in T_{j,k} \mid X=x)$ denote the conditional probabilities of the baseline and alternative-policy level sets. The following overlap assumption bounds the inverse probability weights in the debiased terms.

\noindent\textit{Remark (notation).} $\pi_k(x)=\mathbb P_{\mathrm{obs}}(Z\in S_k\mid X=x)$ is a population conditional probability taking values in $[0,1]$, with fitted baseline nuisances indexed by $k=1,\ldots,K$; no $\pi_0$ regression or inverse-probability correction is fitted. In Lemma~\ref{lem:orthogonality_discrete}, Theorem~\ref{thm:dml_normality}: $\pi_k(x)\geqslant c_\pi>0$ for every $x\in\mathcal X$ and $k=1,\ldots,K$. $\tau_{j,k}(x)=\mathbb P_{\mathrm{obs}}(Z\in T_{j,k}\mid X=x)$ is a population conditional probability in $[0,1]$, with fitted policy-set nuisances indexed by nonempty $T_{j,k}$, $k=0,\ldots,K$ and $j=1,\ldots,l_k$; structural boundary sets are empty and have probability zero. In Lemma~\ref{lem:orthogonality_discrete}, Theorem~\ref{thm:dml_normality}: $\tau_{j,k}(x)\geqslant c_\pi>0$ for every $x\in\mathcal X$ and every nonempty $T_{j,k}$.


\begin{assumption}[Overlap]\label{asp:overlap}
    There exists a constant $c_\pi > 0$ such that $\pi_k(x) \geqslant c_\pi$ for every $x \in \mathcal{X}$ and observed positive-propensity level $k \in \{1, \dots, K\}$, and $\tau_{j,k}(x)\geqslant c_\pi$ for every nonempty alternative-policy level set $T_{j,k}$.
\end{assumption}

\noindent\textit{Remark.} For every $x\in\mathcal X$ and every nonempty $T_{j,k}$, $\tau_{j,k}(x)\geqslant c_\pi$, where $c_\pi>0$ is the same constant used for $\pi_k$ in \cref{asp:overlap}.

\noindent\textit{Remark (notation).} $c_\pi$ is a common positive overlap constant for the specified nonempty level sets, with no overlap requirement on the structural index $k=0$. In Theorem~\ref{thm:dml_normality}: The same $c_\pi$ determines the clipping interval $[c_\pi/2,1]$ for $\hat\pi_k$ and $\hat\tau_{j,k}$.

The index $k=0$ is the structural convention $p_0\equiv0$ and need not correspond to a nonempty instrument level set. Accordingly, every minus-side conditional expectation, inverse-probability correction, and propensity residual indexed by $k=0$ is omitted from the scores; it is multiplied by the known boundary mass $p_0=0$.

\subsubsection{Discrete Outcome}\label{subsubsec:discrete_outcome}

When the outcome $Y$ takes finitely many values---in particular, when $Y \in \{0,1\}$ is binary---the conditional quantile function $Q_{Y,k\mid X}$ becomes a step function with atoms, and the continuity hypothesis of \cref{prop:J_specialization}(i) fails. Instead, \cref{prop:J_specialization}(ii) applies: the representation~(\ref{eq:binary_lower_bound}) expresses $J_{\text{full},k}$ as a difference of observed joint probabilities $P_{1,k+1} - P_{1,k}$ and evaluates the sub-interval contribution in closed form through $h_{j,k}^{-}$, so the nuisance list collapses to conditional expectations alone and quantile estimation is avoided entirely. We present the binary case; the extension to general discrete outcomes is analogous.

The estimator is again~(\ref{eq:estimator_dml}), now with the scores of \cref{app:scores_discrete}. The nuisance-estimation procedure mirrors \cref{alg:dml_continuous_outcome} with the conditional-quantile and indicator-expectation steps removed, and with one addition: because the $\hat P_{1,k}$ are fitted on different instrument subsets, they need not be monotone in $k$ in finite samples, so they are clipped to $[0,1]$ and rearranged before $\hat J_{\mathrm{full},k}=\hat P^*_{1,k+1}-\hat P^*_{1,k}$ is formed. \cref{alg:dml_binary_outcome} states the procedure and \cref{app:procedure_binary_outcome} and \cref{app:scores_discrete} give the full details and the orthogonal scores.

\noindent\textit{Remark (notation).} $\hat P=\hat p(Z,X)$ denotes the fitted baseline propensity index, with observation-indexed notation $\hat P_i=\hat p(Z_i,X_i)$ and values in $[0,1]$. In Theorem~\ref{thm:continuous_rate}, Corollary~\ref{cor:continuous_primitive_rate}: The fitted map $\hat p$ is held fixed conditional on the training samples $I_0,I_1$ when evaluating expressions involving $\hat P=\hat p(Z,X)$ under the target distribution of $(Z,X)$. In Corollary~\ref{cor:continuous_primitive_rate}: The fitted map $\hat p$ is trained on $I_0$, so the selection rule involving $\hat P$ is fixed conditional on $I_0$.


To ensure differentiability of the target functional with respect to the nuisance parameters, we impose the following gap condition, which replaces the conditional-density smoothness required by part~\textup{(i)} of \cref{thm:dml_normality}.

\begin{assumption}[Lower-Boundary Discrete Outcome Gap]\label{asp:discrete_outcome_gap}
    There exists $c_h>0$ such that, for every $k=0,\ldots,K-1$ with $l_k\geqslant1$ and $j=0,\ldots,l_k$, almost surely,
    \[
    \left|p_{k+1}(X)-J_{\mathrm{full},k}(X)-q_{j,k}(X)\right|\geqslant c_h.
    \]
\end{assumption}

\noindent\textit{Remark (notation).} $c_h$ is one positive constant uniform over the indicated indices and almost every covariate value, including the fixed lower boundary $q_{0,k}=p_k$.

This separation makes branch-selection errors for the lower bound second order under the stated $L_2$ rates, including at the fixed lower boundary $q_{0,k}=p_k$. It guarantees that the lower target is pathwise-differentiable in the nuisances despite the $\max$ operators appearing in $h_{j,k}^{-}$. The symmetric upper estimator, when used, requires the analogous separation $|p_k+J_{\mathrm{full},k}-q_{j,k}|\geqslant c_h$ for $j=1,\ldots,l_k+1$.

\subsubsection{Asymptotic Normality}\label{subsubsec:dml_normality}

The two regimes share a common set of conditions---the structural assumptions, the level-set separation and overlap conditions, and $n^{-1/4}$ nuisance rates with clipping---and differ only in what is required at the sub-interval boundaries: a conditional density near each boundary quantile when the outcome is continuous, and the gap condition \cref{asp:discrete_outcome_gap} when it is discrete. A single theorem therefore covers both; both parts are proved in \cref{app:proofs_est_inf}.

\begin{theorem}[Asymptotic Normality] \label{thm:dml_normality}
    Suppose \cref{asp:structure}, \cref{asp:phi_smooth}, \cref{asp:prop_gap}, and \cref{asp:overlap} hold and $\mathcal{Z}$ is finite. Suppose the pre-recovery first stage fitted on $I_1$ satisfies $\max_{z\in\mathcal Z}\|\hat p(z,\cdot)-p(z,\cdot)\|_{L_2}=o_P(n^{-1/4})$. Given the recovered level sets, suppose the remaining score nuisances fitted on $I_1$ are clipped to their known bounded ranges, every estimated conditional probability is clipped to $[0,1]$ before any difference is formed, $\hat\pi_k$ and $\hat\tau_{j,k}$ are clipped further to $[c_\pi/2,1]$, and the estimators converge componentwise at rate $o_P(n^{-1/4})$ in $L_2$. Suppose in addition that one of the following holds.
    \begin{enumerate}[label=\textup{(\roman*)}]
        \item \textbf{Continuous outcome.} The nuisances include the boundary quantiles $\nu_{j,k}$, and, uniformly in $X$, each relevant conditional law $Y\mid X,Z\in S_{k'},W=1$ has a density bounded in a fixed neighborhood of every true boundary $\nu_{j,k}(X)$, $j=1,\ldots,l_k$, for $k'=k+1$ and also for $k'=k$ when $k\geqslant1$; at the fixed boundary $\nu_{0,k}=y_{\min}$, these conditional laws have no atom at $y_{\min}$.
        \item \textbf{Discrete outcome.} The outcome is binary, \cref{asp:discrete_outcome_gap} holds, and the nuisances are $p_k$, $P_{1,k}$, $\pi_k$, $\tau_{j,k}$, $a_{j,k}$, $\gamma_{\mathrm{full},k}$, $\gamma_{j,k}$, and $\gamma_K$.
    \end{enumerate}
    Then the sample-split estimator (\ref{eq:estimator_dml}) is $\sqrt{n}$-consistent and asymptotically normal:
    \begin{equation*}
        \sqrt{n/2} \big( \hat{\underline{\theta}}_{\omega, 1} - \underline{\theta}_{\omega, 1} \big) \xrightarrow{d} \mathcal{N}(0, \sigma^2),
    \end{equation*}
    where the asymptotic variance is the variance of the true orthogonal score, $\sigma^2 = \mathbb{E}\big[ ({\psi}(O; \eta) - \underline{\theta}_{\omega, 1})^2 \big]$. It can be consistently estimated by its sample analog $\hat{\sigma}^2 = \frac{1}{|I_2|} \sum_{i \in I_2} \big(\hat{\psi}(O_i; \hat{\eta}) - \hat{\underline{\theta}}_{\omega, 1}\big)^2$.

\emph{Additional condition.} $K\geqslant1$.
\end{theorem}

\noindent\textit{Remark.} Every nuisance-error norm in this theorem is the population norm $\|\hat f-f\|_{L_2}^2=\mathbb E_X[|\hat f(X)-f(X)|^2]$, with the fitted function held fixed during integration. For a fresh $X$ with the population marginal law and independent of $I_1$, $\|\hat p(z,\cdot)-p(z,\cdot)\|_{L_2}^2=\mathbb E[(\hat p(z,X)-p(z,X))^2\mid I_1]$, and $n^{1/4}\max_{z\in\mathcal Z}\|\hat p(z,\cdot)-p(z,\cdot)\|_{L_2}\xrightarrow{P}0$ as $n\to\infty$. Pointwise in $x$, the fitted probability, propensity, share, and indicator-mean nuisances lie in $[0,1]$, $\hat\gamma_{\mathrm{full},k}(x)\in[-1,1]$, $\hat\nu_{j,k}(x),\hat J_{\mathrm{full},k}^{\pm}(x)\in[y_{\min},y_{\max}]$, and $\hat J_{j,k}^{\pm}(x)\in[\min\{0,y_{\min}\},\max\{0,y_{\max}\}]$, for each applicable fitted component; denominator clipping is applied further as stated. Componentwise $L_2$ errors integrate over the marginal law of $X$ under $\mathbb P_{\mathrm{obs}}$, holding each fitted function and any instrument argument $z$ fixed; $o_P(n^{-1/4})$ refers to convergence in probability over the observed sample and random partition as $n\to\infty$. $Y\in\{0,1\}$ almost surely. In this definition, $\psi_{\mathrm{full},k}^{\mathrm{prod}}$ and $\psi_{j,k}^{\mathrm{prod}}$ are the binary-outcome components of Section~\ref{app:scores_discrete}, and $\psi_{\Delta_K}$ is the common tail component of Section~\ref{app:scores_continuous}. The binary nuisance list supplements the separately fitted first-stage propensity, so $\eta(x)=\{p(z,x),p_k(x),P_{1,k}(x),\pi_k(x),\tau_{j,k}(x),a_{j,k}(z,x),\gamma_{\mathrm{full},k}(x),\gamma_{j,k}(x),\gamma_K(x)\}_{j,k}$, with $p_0=P_{1,0}=0$ fixed and no fitted $\pi_0$. The estimated baseline bins used to index $\widehat T_{j,k}$ include their left endpoints and exclude their right endpoints for $k<K$; the final bin for $k=K$ includes both endpoints, including $1$. Two remaining labeled vertices are connected exactly when their averaged values differ by less than $c_{\mathrm{gap}}/2$. The structural-zero assignment in Algorithm~\ref{alg:dml_continuous_outcome} is $\widehat S_0=\{z:\bar p_2(z)\leqslant\rho_n\}$, with deterministic $\rho_n\downarrow0$ satisfying $n^{1/4}\rho_n\to\infty$. At each evaluation covariate $X_i$, $(\hat P^*_{1,1}(X_i),\ldots,\hat P^*_{1,K}(X_i))$ is the nondecreasing sort of the clipped vector $(\hat P_{1,1}(X_i),\ldots,\hat P_{1,K}(X_i))$, with $\hat P^*_{1,0}(X_i)=0$; this same vector is used in every score plug-in and residual, without a separate gap projection. In this definition, all components of $\psi(O;\eta)$ are those of Section~\ref{app:scores_continuous}, in particular $\psi_{j,k}^{\mathrm{prod}}=\gamma_{j,k}\psi_{j-1,k}^{J}+J_{j-1,k}(\indicator(Z\in\bigcup_{l=j}^{l_k}T_{l,k})-\gamma_{j,k})$.


\subsection{Continuous Instrument Setting}
We now consider a connected instrument domain $\mathcal Z$. Write $P=p(Z,X)$ and define
\begin{gather*}
    \underline{p}(x)\coloneqq\inf_{z\in\mathcal Z}p(z,x),
    \qquad
    \overline{p}(x)\coloneqq\sup_{z\in\mathcal Z}p(z,x), \\
    q^{\mathrm{tr}}(z,x)\coloneqq\min\{\overline{p}(x),\max\{\underline{p}(x),q(z,x)\}\}.
\end{gather*}
The standing compactness of $\mathcal Z$ and continuity of $z\mapsto p(z,x)$ imply that the conditional support of $P$ is the interval $[\underline{p}(x),\overline{p}(x)]$. Define $g_1(u,x)\coloneqq g_1^{\mathrm c}(u,x)=\int_{\mathcal Y}y\,H_{1,x}(u,\mathrm dy)$, the canonical absolutely continuous version of $\mathbb E_{\mathrm{obs}}[YW\mid P=u,X=x]$. The truncated policy propensity $q^{\mathrm{tr}}$ is needed because this function is identified only on the conditional support. For $\underline{p}(x)>0$, let
\begin{equation}\label{eq:boundary_treated_law}
    \mu_{1,\underline p\mid x}(\mathrm dy)
    \coloneqq
    \frac{H_{1,x}(\underline{p}(x),\mathrm dy)}{\underline{p}(x)},
\end{equation}
and denote its quantile function by $Q_{1,\underline p\mid x}$. Equivalently, it is the canonical boundary limit of $\mathcal L(Y\mid W=1,P=p,X=x)$ as $p\downarrow\underline p(x)$ through the conditional support: it averages the law of $Y(1)$ over $0<U\leqslant\underline p(x)$. When $\underline{p}(x)=0$, choose $\mu_{1,\underline p\mid x}$ arbitrarily; every expression below multiplies its contribution by $\underline{p}(x)$ and is therefore invariant to this choice.

The general bound~(\ref{eq:covariate_continuous_lower_bound}) contains the derivative $\partial_u g_1(u,X)$. The next proposition integrates this derivative analytically, so the estimator never differentiates a nonparametric regression.

\begin{proposition}[Continuous-Instrument Reduction via Fubini]\label{prop:continuous_reduction}
    Suppose \cref{asp:structure} and \cref{asp:continuity_p} hold, and $\mathcal{Z}$ is connected. Then the sharp lower bound in \cref{thm:covariate_continuous_lower_bound} admits the representation
    \begin{equation} \label{eq:continuous_simplified_target}
     \begin{split}
            \underline{\theta}_{\omega, 1} &= \mathbb{E}_{Z,X} \Bigg[ - \int_{\min\{\underline{p}(X), q(Z,X)\}}^{\underline{p}(X)} Q_{1,\underline p\mid X}\!\left({v}/{\underline{p}(X)}\right) \, \mathrm{d}v \\
        &\qquad + g_1\big(q^{\mathrm{tr}}(Z,X), X\big) - g_1\big(P, X\big)  \\
        &\qquad + y_{\min} \max\{0, q(Z,X) - \overline{p}(X)\} \Bigg],
     \end{split}
    \end{equation}
    where the first integral is defined as zero when $\underline{p}(X)=0$. The sharp upper bound is obtained by replacing $Q_{1,\underline p\mid X}(v/\underline{p}(X))$ with $Q_{1,\underline p\mid X}(1-v/\underline{p}(X))$ and $y_{\min}$ with $y_{\max}$. 
\end{proposition}

\noindent\textit{Remark.} For $\underline p(x)>0$ and $0<t<1$, $Q_{1,\underline p\mid x}(t)=\inf\{y\in\mathbb R:\mu_{1,\underline p\mid x}(\mathcal Y\cap(-\infty,y])\geqslant t\}$; at $t=0,1$ use the endpoint extension of the initial-gap outcome quantile in the referenced bound. For $0<t<1$, use the standard lower generalized quantile $Q_{1,\underline p\mid x}(t)=\inf\{y\in\mathbb R:\mu_{1,\underline p\mid x}(\mathcal Y\cap(-\infty,y])\geqslant t\}$.


The first term uses the treated boundary law in (\ref{eq:boundary_treated_law}); the second evaluates $g_1$ only at $P$ and at the support projection $q^{\mathrm{tr}}(Z,X)$. This lower truncation matters for policies with $q(Z,X)<\underline{p}(X)$: the movement below $\underline{p}(X)$ belongs to the OT gap term, whereas the LIV contribution stops at $\underline{p}(X)$. The third term covers $q(Z,X)>\overline{p}(X)$. The proof is given in \cref{app:procedure_fubini} of \oalocation.

\cref{prop:continuous_reduction} motivates a plug-in estimator, obtained by replacing each nuisance in~(\ref{eq:continuous_simplified_target}) with a machine-learning estimate trained on auxiliary samples. We use three-way sample splitting with $|I_j|/n\to\rho_j\in(0,1)$ for $j=0,1,2$: the sample $I_0$ estimates $p$, $q$, $\underline{p}$, and $\overline{p}$; conditional on $I_0$, the sample $I_1$ estimates $g_1$ and the treated boundary law; and $I_2$ evaluates the target. The estimator is
\begin{equation}\label{eq:estimator_continuous_instrument}
    \hat{\underline{\theta}}_{\omega, 1} = \frac{1}{|I_2|} \sum_{i \in I_2} \Big( \hat{\psi}_{1}(O_i) + \hat{\psi}_{2}(O_i) + \hat{\psi}_{3}(O_i) \Big),
\end{equation}
where $\hat\psi_1,\hat\psi_2,\hat\psi_3$ substitute estimated nuisances into the three terms of (\ref{eq:continuous_simplified_target}). Throughout, $\hat P=\hat p(Z,X)$ and $\hat q=\phi(Z,X,\hat p(Z,X))$, while $\hat{\underline p}$, $\hat{\overline p}$, $\hat q^{\mathrm{tr}}$, $\min\{\hat{\underline p}(X),\hat q(Z,X)\}$, and $(\hat q(Z,X)-\hat{\overline p}(X))_+$ are obtained by substituting these estimates into the corresponding definitions. We clip estimated propensity scores and nuisance predictions to their known supports.

\noindent\textit{Remark (notation).} $\hat{\underline p}(x)=\inf_{z\in\mathcal Z}\hat p(z,x)$ is the fitted lower propensity endpoint, a data-dependent function of $x$ taking values in $[0,1]$, held fixed conditional on the training samples $I_0,I_1$. In Corollary~\ref{cor:continuous_primitive_rate}: $\hat{\underline p}$ is obtained from the propensity fit trained on $I_0$ and is therefore fixed conditional on $I_0$ in the selected-law definition. $\hat{\overline p}(x)=\sup_{z\in\mathcal Z}\hat p(z,x)\in[0,1]$ is the fitted upper propensity endpoint, with the fitted function held fixed in target evaluation conditional on $I_0,I_1$. In Corollary~\ref{cor:continuous_primitive_rate}: The endpoint is computed from $\hat p$ fitted on $I_0$. $\hat q^{\mathrm{tr}}(z,x)=\min\{\hat{\overline p}(x),\max\{\hat{\underline p}(x),\hat q(z,x)\}\}$ is the fitted policy propensity projected onto $[\hat{\underline p}(x),\hat{\overline p}(x)]$, with its fitted map held fixed in target evaluation conditional on $I_0,I_1$. In Corollary~\ref{cor:continuous_primitive_rate}: The fitted map is constructed on $I_0$ from the clipped $\hat p$ and $\hat q(z,x)=\phi(z,x,\hat p(z,x))$. $\hat\psi_1(O)$ denotes the estimated lower-bound boundary contribution $-\int_{\min\{\hat{\underline p}(X),\hat q(Z,X)\}}^{\hat{\underline p}(X)}\widehat Q_{1,\underline p\mid X}(v/\hat{\underline p}(X))\,\mathrm dv$, with value zero when $\hat{\underline p}(X)=0$. In Corollary~\ref{cor:continuous_primitive_rate}: For $\hat{\underline p}(X_i)>0$, the concrete implementation evaluates $-\hat{\underline p}(X_i)[\widehat C(1\mid X_i)-\widehat C(\hat\tau_i\mid X_i)]$, where $\hat\tau_i=\min\{1,\hat q(Z_i,X_i)/\hat{\underline p}(X_i)\}$; cumulative-integral grid values are computed by the trapezoid rule, and off-grid lookup uses linear interpolation between adjacent cumulative-integral grid values. For $\hat{\underline p}(X_i)=0$, it returns zero without evaluating $\hat\tau_i$. $\hat\psi_3(O)=y_{\min}(\hat q(Z,X)-\hat{\overline p}(X))_+$ is the fitted lower-bound tail score, evaluated at $O_i$ for $i\in I_2$ with the fitted score map held fixed conditional on $I_0,I_1$. For $j\in\{0,1,2\}$, $\rho_j\in(0,1)$ is the limiting split proportion $|I_j|/n\to\rho_j$, with $\rho_0+\rho_1+\rho_2=1$.


We state the convergence rate under high-level nuisance-rate assumptions: (\ref{eq:estimator_continuous_instrument}) is well defined for any nuisance estimators trained on $I_0\cup I_1$. \cref{alg:continuous_instrument} then provides one concrete construction, and \cref{cor:continuous_primitive_rate} establishes its rate under conditions on those particular learners, collected in \cref{app:continuous_learner_conditions} of \oalocation.
\begin{assumption}[Targeted Nuisance Rates]\label{asp:continuous_rates}
    Conditional on the training samples $I_0\cup I_1$, suppose:
    \begin{enumerate}
        \item Each of the three maps entering the boundary integral and the trivial tail,
        \begin{equation*}
            \Lambda \;\in\; \Bigl\{\ \underline{p}(X),\ \ \min\{\underline{p}(X),\,q(Z,X)\},\ \ \bigl(q(Z,X)-\overline{p}(X)\bigr)_+ \ \Bigr\},
        \end{equation*}
        satisfies
        \[
            \mathbb E_{Z,X}\!\left[\,\bigl|\hat\Lambda-\Lambda\bigr|\ \middle|\ I_0,I_1\right]=O_P(r_{e,n}),
        \]
        with $r_{e,n}\to0$.
        \item The clipped estimator $\hat g_1$ satisfies the target-evaluation rate
        \begin{multline*}
            \mathbb E_{Z,X}\Bigl[
              |\hat g_1(\hat q^{\mathrm{tr}}(Z,X),X)-g_1(q^{\mathrm{tr}}(Z,X),X)| \\
              +|\hat g_1(\hat P,X)-g_1(P,X)|
              \ \Bigm|\ I_0,I_1\Bigr]=O_P(r_{g,n}).
        \end{multline*}
        Here $r_{g,n}\to0$.
        \item Let $\hat\mu_{1,\underline p\mid x}$ be the estimated treated boundary law with quantile function $\widehat Q_{1,\underline p\mid x}$. It satisfies the target-weighted rate
        \[
            \mathbb E_X\!\left[\underline{p}(X)W_1\!\left(\hat\mu_{1,\underline p\mid X},\mu_{1,\underline p\mid X}\right)\ \middle|\ I_0,I_1\right]=O_P(r_{\mu,n}).
        \]
        Here $r_{\mu,n}\to0$.
    \end{enumerate}
\end{assumption}

\noindent\textit{Remark (notation).} $\hat g_1(u,x)$ is an estimated regression on $[0,1]\times\mathcal X$, clipped to the known range $[\min\{0,y_{\min}\},\max\{0,y_{\max}\}]$ of $YW$, and evaluated at the estimated target indices $\hat q^{\mathrm{tr}}(Z,X)$ and $\hat P$. In Corollary~\ref{cor:continuous_primitive_rate}: At $u_m=m/M_g$, $m=0,\ldots,M_g$, the fitted value is the intercept in $(\hat g_1(u_m,\cdot),\hat\beta_m)\in\arg\min_{g,\beta\in\mathcal F_n}\sum_{i\in I_1}K_{h_n}(\hat P_i-u_m)[Y_iW_i-g(X_i)-\beta(X_i)(\hat P_i-u_m)]^2$; between grid points it is the clipped linear interpolant. The learner loss compares $\hat g_1(\hat U_\ell,X)$ with $g_1(\Pi_X(\hat U_\ell),X)$ for $\ell\in\{1,2\}$. $\hat\mu_{1,\underline p\mid x}$ is an estimated probability law on $\mathcal Y$, indexed by $x\in\mathcal X$, whose quantile function is $\widehat Q_{1,\underline p\mid x}$ and whose target is the canonical treated boundary law; its target-error loss is weighted by $\underline p(x)$, so its value where $\underline p(x)=0$ does not affect the target. In Corollary~\ref{cor:continuous_primitive_rate}: The law is generated by conditional check-function fits on $\mathcal I_{\delta,1}=\{i\in I_1:W_i=1,\hat P_i\leqslant\hat{\underline p}(X_i)+\delta_n\}$ at $\tau_m=m/M_Q$, $m=1,\ldots,M_Q-1$, clipped to $[y_{\min},y_{\max}]$, sorted at each $x$, and linearly interpolated between adjacent quantile levels; endpoint values may be $y_{\min},y_{\max}$ or the smallest and largest fitted values. $\widehat Q_{1,\underline p\mid x}$ denotes the quantile function of $\hat\mu_{1,\underline p\mid x}$, with values in $[y_{\min},y_{\max}]$ and endpoint values supplied by the chosen endpoint convention. In Corollary~\ref{cor:continuous_primitive_rate}: At $\tau_m=m/M_Q$, $m=1,\ldots,M_Q-1$, fit $\widetilde Q(\tau_m\mid\cdot)\in\arg\min_{f\in\mathcal F_Q}\sum_{i\in\mathcal I_{\delta,1}}\rho_{\tau_m}(Y_i-f(X_i))$, setting every grid fit to $y_{\min}$ if the selected sample is empty; clip to $[y_{\min},y_{\max}]$, sort the fitted values for each $x$, and interpolate linearly in $\tau$, with endpoints either $y_{\min},y_{\max}$ or the smallest and largest fitted values. $r_{e,n}$ is a nonnegative rate sequence tending to zero controlling each of the three training-conditional absolute-error losses for the boundary and tail maps. In Corollary~\ref{cor:continuous_primitive_rate}: $r_{e,n}=O(r_{p,n})$ for the concrete nuisance construction. $r_{g,n}$ is a nonnegative rate sequence tending to zero for the sum of the two training-conditional target-evaluation losses, comparing fitted regressions at estimated indices with population regressions at the corresponding true indices. In Corollary~\ref{cor:continuous_primitive_rate}: $r_{g,n}=O(r_{g,\mathrm{loc},n}+r_{p,n})$ for the concrete nuisance construction. $r_{\mu,n}$ is a nonnegative deterministic rate sequence tending to zero for the training-conditional, $\underline p(X)$-weighted Wasserstein loss against the canonical treated boundary law. In Corollary~\ref{cor:continuous_primitive_rate}: For the concrete nuisance construction, one may take $r_{\mu,n}=O(r_{Q,n}+\delta_n+r_{p,n}+\varepsilon_n+M_Q^{-1})$. $\Lambda$ is population notation whose meaning is determined by its local introduction: an evaluated endpoint or tail expression, or the logistic function in the simulation index. In Theorem~\ref{thm:continuous_rate}: $\Lambda$ ranges separately over $\underline p(X)$, $\min\{\underline p(X),q(Z,X)\}$, and $(q(Z,X)-\overline p(X))_+$, with the error condition required for each choice; these are evaluated population expressions under the law of $(Z,X)$. $\hat\Lambda$ denotes the corresponding estimated endpoint or tail expression evaluated at $(Z,X)$; in $\mathbb E_{Z,X}[|\hat\Lambda-\Lambda|\mid I_0,I_1]$, the fitted functions are held fixed and the expectation integrates the evaluation law of $(Z,X)$.


Part~(1) controls the two limits of the boundary integral and the positive-part tail, averaged under the target distribution. Part~(2) addresses the more demanding generated-regressor problem. The estimated propensity score plays two roles: $\hat P$ is used in place of the latent index $P$ when learning $g_1$, and it determines the random arguments at which $\hat g_1$ is evaluated. Part~(3) controls the boundary OT contribution in its natural loss. Indeed, in one dimension $W_1$ equals the integrated absolute quantile error, so this assumption essentially requires that the estimated boundary quantile function $\widehat Q_{1,\underline p\mid x}$ converge in $L_1$ to the true boundary quantile function $Q_{1,\underline p\mid x}$, weighted by the target $\underline{p}(X)$.

\begin{theorem}[Nonparametric Convergence Rate]\label{thm:continuous_rate}
    Suppose \cref{asp:structure}, \cref{asp:continuity_p}, and \cref{asp:continuous_rates} hold. Then
    \[
        |\hat{\underline\theta}_{\omega,1}-\underline\theta_{\omega,1}|
        =O_P\!\left(n^{-1/2}+r_{e,n}+r_{g,n}+r_{\mu,n}\right).
    \]
\end{theorem}

\noindent\textit{Remark.} For every probability measure $\mu$ on $\mathbb R$ and $0<t<1$, $Q_\mu(t)=\inf\{y\in\mathbb R:\mu(( -\infty,y])\geqslant t\}$; the integral defining $W_1$ is a Lebesgue integral over $(0,1)$ and requires no endpoint values.


\cref{thm:continuous_rate} provides a high-level theoretical guarantee, and \cref{app:procedure_continuous_instrument} gives a concrete construction: on $I_0$ we regress $W$ on $(Z,X)$; on $I_1$ we estimate $g_1$ by boundary-adaptive local-linear smoothing in the scalar index $P$ combined with a flexible regression in $X$, and we estimate the treated boundary law by conditional quantile regression on the treated observations whose estimated propensity lies within $\delta_n$ of the estimated lower boundary. Under conditions (L1)--(L3), \cref{cor:continuous_primitive_rate} relates its learner errors to the target error:
\[
    |\hat{\underline\theta}_{\omega,1}-\underline\theta_{\omega,1}|
    =O_P\!\left(n^{-1/2}+r_{p,n}+r_{g,\mathrm{loc},n}
    +r_{Q,n}+\delta_n+\varepsilon_n+M_Q^{-1}\right),
\]
where $r_{p,n}$ is a sup-norm first-stage rate, $r_{g,\mathrm{loc},n}$ is an $L_1$ prediction rate for $\hat g_1$ at the two arguments at which the target evaluates it, and $r_{Q,n}$ is a $\underline p(X)$-weighted $L_1$ rate for the conditional quantile fit relative to the law selected by the estimated boundary window. The threshold-crossing model gives the linear localization error $\delta_n+r_{p,n}$ directly; $\varepsilon_n$ and $M_Q^{-1}$ are the trimmed quantile-tail and grid-interpolation errors.

\noindent\textit{Remark (notation).} The notation $\delta$ denotes a nuisance perturbation direction, whereas the indexed notation $\delta_n$ denotes a fixed boundary-localization window sequence. In Lemma~\ref{lem:orthogonality_discrete}, Theorem~\ref{thm:dml_normality}: $\delta$ lies in the same function space as $\eta$ and satisfies that $\eta_0+t\delta$ is admissible for every $t$ near zero; fixed structural components have zero perturbation, and within-set share components satisfy $\sum_{z\in T_{j,k}}\delta a_z(X)=0$. In Corollary~\ref{cor:continuous_primitive_rate}: $\delta_n\to0$ is the localization window in $\mathcal I_{\delta,1}=\{i\in I_1:W_i=1,\ \hat P_i\leqslant\hat{\underline p}(X_i)+\delta_n\}$, indexed by sample size $n$. $r_{p,n}$ is a nonnegative deterministic sample-size rate for uniform propensity estimation, rather than an integrated prediction loss. In Corollary~\ref{cor:continuous_primitive_rate}: $r_{p,n}\to0$. $r_{g,\mathrm{loc},n}$ is a nonnegative deterministic rate sequence indexed by sample size for prediction under both target evaluation-index laws, using the projected estimated index in the population comparator. In Corollary~\ref{cor:continuous_primitive_rate}: $r_{g,\mathrm{loc},n}\to0$. $r_{Q,n}$ is a nonnegative deterministic rate sequence indexed by sample size, measuring the interior-grid quantile-fit loss relative to the selected-law target, rather than the canonical boundary law. In Corollary~\ref{cor:continuous_primitive_rate}: $r_{Q,n}\to0$. $\delta_n>0$ is the fixed, sample-size-indexed localization window used in the fitted selection rule $\hat P\le\hat{\underline p}(X)+\delta_n$. In Corollary~\ref{cor:continuous_primitive_rate}: $\delta_n\to0$ as $n\to\infty$. $\varepsilon_n$ is the quantile-analysis trimming sequence satisfying $M_Q^{-1}\le\varepsilon_n<1/2$ and $\varepsilon_n\to0$; it restricts the learner analysis to intermediate quantiles without changing the estimator. $M_Q$ is the integer resolution of the quantile grid $\tau_m=m/M_Q$, $m=1,\ldots,M_Q-1$, and may vary with $n$. In Corollary~\ref{cor:continuous_primitive_rate}: $M_Q^{-1}\to0$, so the quantile grid resolution diverges with sample size.


Applying the same construction to the reversed system $W^0=1-W$, $P^0=1-P$, $q^0=1-q$, $U^0=1-U$ gives the untreated component, and the two bounds for the full PRTE converge at the maximum of the two componentwise rates; see \cref{app:procedure_continuous_instrument}.
